\begin{textsample}{NEG Dim 1 – 2014\_c5 – Score 8.00 – 2014\_c5/enade\_2014\_eng\_alim\_1adfbde86f.txt}  \label{ex:f1_neg_339}
CAPÍTULO 5 DISTRIBUIÇÃO DOS CONCEITOS 5.1 PANORAMA NACIONAL DA DISTRIBUIÇÃO DOS CONCEITOS O Conceito Enade é um indicador de qualidade que avalia o desempenho dos estudantes a partir dos resultados do ENADE.
Os valores possíveis vão de 1 ( pior situação ) a 5 ( melhor situação ).
As Unidades de Observação com apenas um ou sem nenhum concluinte participante não obtêm o Conceito Enade, ficando Sem Conceito ( SC ).
A Tabela 5.1 apresenta a quantidade e distribuição de cursos de Engenharia de Alimentos participantes do ENADE/2014, por faixa de conceito e Grande Região.
Mais uma vez, ressaltando que a diferença entre os cursos tabulados neste capítulo e no capítulo 2 são os cursos sem conceito, em princípio, aqueles sem alunos concluintes que participassem da prova ou que tivessem tão somente um aluno concluinte21.
Observando -se os dados da Tabela 5.1, nota -se que, dos 67 cursos participantes, 22 ( 32,8 \% ) classificaram -se com conceito 3, o valor modal.
Este foi também o conceito modal em mais duas regiões : Nordeste ( 60,0 \% ) e Sul ( 36,4 \% ).
Nas regiões Norte e Centro-Oeste, a moda foi o conceito 2, e na região Sudeste, o conceito 4.
Em nível nacional, o conceito 2 foi o segundo mais frequente ( 28,4 \%, correspondendo a 19 cursos ) e o conceito 4, o terceiro ( 23,9 \%, correspondendo a 16 cursos ).
Houve, ainda, cinco cursos ( 7,5 \% ) que receberam conceito 1 e outros cinco cursos que receberam conceito 5.
Dos 67 cursos de Engenharia de Alimentos, nenhum ficou sem conceito ( SC ). 21 Estes cursos com somente um concluinte não foram considerados no capítulo 2.
A região Norte participou com três cursos ou 4,5 \% do total nacional.
Nesta região, o conceito 2 foi atribuído a dois cursos dos três participantes, o que equivale a 66,7 \% do total regional.
O outro curso recebeu conceito 3.
A região Nordeste participou com dez cursos ou 14,9 \% do total nacional.
Nesta região, seis cursos, 60,0 \% em termos regionais, obtiveram conceito 3, o conceito modal.
O conceito 2 foi atribuído a dois cursos ( 20 \% ), e o conceito 4, a outros dois.
Logo, nenhum curso ficou sem conceito ( SC ) ou recebeu os demais conceitos nessa região.
Os cursos da região Sudeste ficaram distribuídos em todas as faixas de conceitos.
Dos 24 cursos participantes na região ( 35,8 \% do total nacional ), o conceito modal foi 4, obtido por 8 cursos ( 33,3 \% ).
O conceito 1 foi atribuído a dois cursos ( 8,3 \% ), o conceito 2, a cinco cursos ( 20,8 \% ), o conceito 3, a seis cursos ( 25,0 \% ) e o conceito 5, a três cursos ( 12,5 \% ).
Nessa região, nenhum curso ficou sem conceito ( SC ).
Dos 22 cursos da região Sul houve predominância do conceito 3 com 36,4 \% ( conceito modal ), ou seja, oito dos cursos participantes.
Apenas um curso ( 4,5 \% ) recebeu conceito 1.
Os conceitos 2, 4 e 5 foram atribuídos a, respectivamente, cinco, seis e dois cursos, correspondendo a 22,7 \%, 27,3 \% e 9,1 \% dos cursos da região.
Nenhum curso na região Sul ficou sem conceito ( SC ).
Dos oito cursos participantes na região Centro-Oeste, cinco ( 62,5 \% ) receberam conceito 2, o conceito modal.
Nesta região, nenhum curso ficou sem conceito ( SC ) ou recebeu os conceitos 4 ou 5.
Foram atribuídos conceito 1 a dois cursos ( 25,0 \% ) e conceito 3 a um curso ( 12,5 \% ).
O Gráfico 5.1 apresenta a distribuição do Conceito Enade segundo Grande Região : poligonais mais à esquerda representam Grandes Regiões com pior distribuição, e poligonais mais à direita, Grandes Regiões com melhores conceitos.
A região Centro-Oeste ( linha preta ) apresenta os cursos em IES com conceitos mais baixos, ao passo que as regiões Sudeste ( linha vermelha ) e Sul ( linha roxa ) apresentam as distribuições com valores semelhantes, e maiores.
Cursos em IES das regiões Norte ( linha azul ) e Nordeste ( linha verde ) apresentam uma situação intermediária. 5.2 CONCEITOS POR CATEGORIA ADMINISTRATIVA E POR GRANDE REGIÃO A Tabela 5.2 apresenta a distribuição dos cursos participantes do ENADE/2014 de Engenharia de Alimentos, por Categoria Administrativa, de acordo com os conceitos por eles alcançados, segundo as Grandes Regiões brasileiras.
Dos 67 cursos participantes, 48 ( 71,6 \% ) eram ministrados em Instituições Públicas e 19 ( 28,4 \% ), em Instituições Privadas.
De acordo com as informações da Tabela 5.2, em termos nacionais, entre Instituições Públicas e Privadas : todas as cinco IES que receberam conceito 5 eram Públicas.
Dos 48 cursos participantes de IES Públicas, o conceito 3 foi o modal, atribuído a 15 cursos ( 31,3 \% ).
Entre os demais cursos de Instituições Públicas participantes, apenas um curso ( 2,1 \% ) recebeu conceito 1, 13 cursos ( 27,1 \% ) foram avaliados com conceito 2, 14 cursos ( 29,2 \% ) obtiveram conceito 4 e, como já informado, cinco cursos ( 10,4 \% ) receberam conceito 5.
Nesta categoria, nenhum curso ficou sem conceito ( SC ).
Na rede privada, o conceito modal foi o conceito 3, atribuído a sete dos 19 cursos da categoria.
Entre os demais cursos participantes, quatro foram avaliados com conceito 1, seis obtiveram conceito 2 e dois receberam conceito 4.
Nesta Categoria Administrativa, nenhum curso ficou sem conceito ( SC ) ou recebeu conceito 5.
Na análise por região, como se registrou no Capítulo 2, na região Norte, todos os três cursos oferecidos eram vinculados a Instituições Públicas.
Desses, dois cursos obtiveram conceito 2, e um curso, conceito 3.
Também na região Nordeste, não houve participação de cursos da rede privada na Área de Engenharia de Alimentos.
Dos dez cursos oferecidos por Instituições Públicas na região, seis ( 60,0 \% ) receberam conceito 3, o valor modal.
Dois cursos ( 20,0 \% ) receberam conceito 2 e outros dois, conceito 4.
Na região Sudeste, a proporção de cursos da rede privada, 41,7 \%, foi mais elevada do que nas demais regiões brasileiras, correspondendo a dez dos 24 cursos participantes.
Nesta categoria, na região Sudeste, os conceitos 2 e 3 foram os conceitos modais, atribuídos, \textbf{cada} qual, a quatro cursos ( 40,0 \% \textbf{cada} ).
Os demais cursos foram avaliados com os conceitos 1 e 4 ( um curso \textbf{cada} ).
Nenhum curso vinculado a IES privada recebeu conceito 5 ou ficou sem conceito ( SC ).
Entre os 14 cursos oferecidos por Instituições Públicas na região Sudeste, o conceito modal foi 4, atribuído a sete cursos.
Os demais cursos receberam os conceitos 1 ( um curso ), 2 ( um curso ), 3 ( dois cursos ) e 5 ( três cursos ).
As Instituições Privadas participaram com sete dos 22 cursos participantes da região Sul, 31,8 \% do total regional.
Desses, três receberam conceito 3, o conceito modal.
Nesta combinação de Categoria Administrativa e Grande Região, um curso obteve conceito 1, dois cursos obtiveram conceito 2, e um curso recebeu conceito 4.
Não foi atribuído conceito 5 a nenhum curso oferecido por IES Privada na região.
As Instituições Públicas na região Sul participaram com 15 cursos ( 68,2 \% do total regional ).
Os conceitos 3 e 4 foram os conceitos modais, \textbf{cada} um dos quais, atribuídos a cinco cursos.
Três cursos foram avaliados com conceito 2, e dois cursos, com conceito 5.
Nenhum curso recebeu conceito 1 ou ficou sem conceito ( SC ).
Na região Centro-Oeste, dois dos oito cursos participantes eram de Instituições Privadas ( 25,0 \% em termos regionais ), ambos avaliados com conceito 1.
Dos seis cursos oferecidos por Instituições Públicas, cinco foram avaliados com conceito 2, conceito modal, e um curso, com conceito 3.
O Gráfico 5.2 apresenta a distribuição do Conceito Enade segundo a Categoria Administrativa da IES.
Os cursos em IES Públicas ( linha azul ) apresentam uma poligonal mais à direita do que os em IES Privadas ( linha verde ) e, consequentemente, uma distribuição de Conceitos Enade com valores maiores. 5.3 CONCEITOS POR ORGANIZAÇÃO ACADÊMICA E POR GRANDE REGIÃO Na Tabela 5.3, encontra -se a distribuição dos conceitos atribuídos aos cursos participantes do ENADE/2014 na Área de Engenharia de Alimentos, por Organização Acadêmica, segundo as Grandes Regiões brasileiras.
Dos 67 cursos de Engenharia de Alimentos participantes, 56 eram oferecidos em Universidades, seis em Centros Universitários e cinco em Faculdades.
Esta distribuição corresponde a, respectivamente, 83,6 \%, 9,0 \% e 7,5 \% dos cursos.
De acordo com os dados apresentados, todos os cinco cursos avaliados com conceito 5 eram vinculados a Universidades.
O conceito modal das Universidades foi 3, atribuído a 18 cursos ( 32,1 \% ).
Os demais cursos de Universidades avaliados receberam os conceitos 1 ( dois cursos ), 2 ( 16 cursos ), 4 ( 15 cursos ) e 5 ( cinco cursos, como já mencionado ).
Entre os seis cursos oferecidos em Centros Universitários, os conceitos modais foram os conceitos 1 e 3, atribuídos, \textbf{cada} qual, a dois cursos ( 33,3 \%, \textbf{cada} ).
Os outros cursos receberam os conceitos 2 e 4 ( um curso, \textbf{cada} ).
Nenhum curso oferecido por Centro Universitário recebeu conceito 5 ou ficou sem conceito ( SC ).
Os cinco cursos de Faculdades participantes ficaram assim avaliados : conceitos 2 e 3 atribuídos, \textbf{cada} qual, a dois cursos e o outro curso restante recebeu conceito 1.
Considerando -se separadamente as regiões brasileiras, verifica -se que, na região Norte, os três cursos participantes eram de Universidades.
Dois desses cursos receberam conceito 2, conceito modal, e um curso obteve conceito 3.
Também na região Nordeste, apenas as Universidades ofereceram cursos na Área de Engenharia de Alimentos.
Dos dez cursos oferecidos, seis obtiveram o conceito modal, 3.
Os demais cursos receberam os conceitos 2 e 4 ( dois cursos \textbf{cada} ).
Na região Sudeste, as Universidades concentraram 15 dos 24 cursos da região.
Entre os cursos em Universidades, nessa região, o conceito modal foi 4, atribuído a sete cursos.
Os demais cursos receberam os conceitos 1 ( um curso ), 2 ( dois cursos ), 3 ( dois cursos ) e 5 ( três cursos ).
Nenhum curso ficou sem conceito ( SC ).
Os Centros Universitários participaram com quatro cursos na região Sudeste, dos quais dois obtiveram o conceito modal, 3.
Os outros dois cursos receberam os conceitos 2 e 4 ( um curso \textbf{cada} ).
As Faculdades foram representadas por cinco cursos na região Sudeste.
Os conceitos 2 e 3 foram os conceitos modais, atribuídos, \textbf{cada} qual, a dois cursos.
O outro curso recebeu conceito 1.
Dos 22 cursos oferecidos por IES na região Sul, 21 eram de Universidades, para os quais o conceito modal foi o conceito 3, atribuído a oito cursos.
Os demais cursos receberam os conceitos 2 ( cinco cursos ), 4 ( seis cursos ) e 5 ( dois cursos ).
Nenhum curso recebeu conceito 1 ou ficou sem conceito ( SC ) nessa categoria e região.
Na região Sul, o único curso na Área de Engenharia de Alimentos não vinculado a Universidades foi oferecido por Centro Universitário, e recebeu conceito 1.
Na região Centro-Oeste, sete dos oito cursos participantes eram oferecidos em Universidades.
O conceito modal foi 2, atribuído a cinco cursos.
Os outros dois cursos receberam conceitos 1 e 3 ( um curso \textbf{cada} ).
Na região o único curso não oferecido por Universidade era de Centro Universitário e obteve conceito 1.
O Gráfico 5.3 apresenta a distribuição do Conceito Enade segundo a Organização Acadêmica da IES.
Os cursos em Universidades ( linha azul ), que foram a maioria dos cursos da Área, apresentam uma poligonal mais à direita do que as demais, denotando uma distribuição de Conceitos Enade maiores.

% matched lemmas: cada
\end{textsample}
